% ============================================================================
% WORKBOOK — Additional: Compound Interest Introduction
% State-specific (TX)
% Practice-focused reinforcement for the study guide topic
% ============================================================================
\section{Compound Interest Introduction}
\topicTitle{Compound Interest Introduction}

\begin{quickReview}{Simple vs.\ Compound Interest}
\textbf{Simple interest} is calculated on the original principal only: $I = Prt$.

\textbf{Compound interest} is calculated on the principal \emph{plus} previously earned interest.
$$A = P(1 + r)^t$$
\begin{itemize}[leftmargin=*]
    \item $A$ = total amount, $P$ = principal, $r$ = annual rate (decimal), $t$ = years
    \item Compound interest grows faster because your interest earns interest!
\end{itemize}

\medskip
\textbf{Example:} $\$500$ at $10\%$ for $3$ years (compounded yearly):

Year 1: $500 \times 1.10 = \$550$ \quad Year 2: $550 \times 1.10 = \$605$ \quad Year 3: $605 \times 1.10 = \$665.50$
\end{quickReview}

% ── Warm-Up ──────────────────────────────────────────────────────────────────
\begin{practiceBox}[funGreen]{\faSun~Warm-Up}

\practiceHeader[funGreen]{Compound Interest — Year by Year}

\prob $\$100$ at $10\%$ compounded yearly. Find the amount after Year 1. \answerBlank[3cm]
\answerExplain{$\$110$}{$100 \times 1.10 = \$110$. Multiply the principal by $1 +$ the rate.}

\prob From Problem 1, find the amount after Year 2. \answerBlank[3cm]
\answerExplain{$\$121$}{$110 \times 1.10 = \$121$. The interest in Year 2 is calculated on $\$110$, not the original $\$100$.}

\prob $\$200$ at $5\%$ compounded yearly. Find the amount after Year 1. \answerBlank[3cm]
\answerExplain{$\$210$}{$200 \times 1.05 = \$210$.}

\prob From Problem 3, find the amount after Year 2. \answerBlank[3cm]
\answerExplain{$\$220.50$}{$210 \times 1.05 = \$220.50$.}

\prob $\$400$ at $8\%$ compounded yearly. Find the amount after Year 1. \answerBlank[3cm]
\answerExplain{$\$432$}{$400 \times 1.08 = \$432$.}
\end{practiceBox}

% ── Practice A: Compound for 2 Years ─────────────────────────────────────────
\begin{practiceBox}{Compound Interest for 2 Years}

\practiceHeader{Calculate year by year. Round to the nearest cent.}

\prob $\$300$ at $6\%$ for $2$ years. \answerBlank[3cm]
\answerExplain{$\$337.08$}{Year 1: $300 \times 1.06 = \$318$. Year 2: $318 \times 1.06 = \$337.08$.}

\prob $\$1{,}000$ at $4\%$ for $2$ years. \answerBlank[3cm]
\answerExplain{$\$1{,}081.60$}{Year 1: $1{,}000 \times 1.04 = \$1{,}040$. Year 2: $1{,}040 \times 1.04 = \$1{,}081.60$.}

\prob $\$750$ at $8\%$ for $2$ years. \answerBlank[3cm]
\answerExplain{$\$874.80$}{Year 1: $750 \times 1.08 = \$810$. Year 2: $810 \times 1.08 = \$874.80$.}

\prob $\$500$ at $12\%$ for $2$ years. \answerBlank[3cm]
\answerExplain{$\$627.20$}{Year 1: $500 \times 1.12 = \$560$. Year 2: $560 \times 1.12 = \$627.20$.}

\prob $\$2{,}000$ at $3\%$ for $2$ years. \answerBlank[3cm]
\answerExplain{$\$2{,}121.80$}{Year 1: $2{,}000 \times 1.03 = \$2{,}060$. Year 2: $2{,}060 \times 1.03 = \$2{,}121.80$.}
\end{practiceBox}

% ── Practice B: Using the Formula ────────────────────────────────────────────
\begin{practiceBox}[funTeal]{Using $A = P(1+r)^t$}

\practiceHeader[funTeal]{Use the compound interest formula. Round to the nearest cent.}

\prob $P = \$600$, $r = 5\%$, $t = 3$ years. Find $A$. \answerBlank[3cm]
\answerExplain{$\$694.58$}{$A = 600(1.05)^3 = 600 \times 1.157625 = \$694.58$.}

\prob $P = \$1{,}000$, $r = 10\%$, $t = 3$ years. Find $A$. \answerBlank[3cm]
\answerExplain{$\$1{,}331$}{$A = 1{,}000(1.10)^3 = 1{,}000 \times 1.331 = \$1{,}331$.}

\prob $P = \$800$, $r = 4\%$, $t = 3$ years. Find $A$. \answerBlank[3cm]
\answerExplain{$\$899.89$}{$A = 800(1.04)^3 = 800 \times 1.124864 = \$899.89$.}

\prob $P = \$1{,}500$, $r = 6\%$, $t = 2$ years. Find $A$. \answerBlank[3cm]
\answerExplain{$\$1{,}685.40$}{$A = 1{,}500(1.06)^2 = 1{,}500 \times 1.1236 = \$1{,}685.40$.}
\end{practiceBox}

% ── Practice C: Simple vs.\ Compound Comparison ─────────────────────────────
\begin{practiceBox}[funOrange]{Simple vs.\ Compound — Side by Side}

\prob $\$500$ at $10\%$ for $2$ years. Find the simple interest total and the compound total.
\answerBlank[5cm]
\answerExplain{Simple: $\$600$; Compound: $\$605$}{Simple: $I = 500 \times 0.10 \times 2 = \$100$; total $= \$600$. Compound: Year 1: $550$; Year 2: $550 \times 1.10 = \$605$.}

\prob $\$1{,}000$ at $8\%$ for $3$ years. Find both totals. \answerBlank[5cm]
\answerExplain{Simple: $\$1{,}240$; Compound: $\$1{,}259.71$}{Simple: $I = 1{,}000 \times 0.08 \times 3 = \$240$; total $= \$1{,}240$. Compound: $A = 1{,}000(1.08)^3 = 1{,}000 \times 1.259712 \approx \$1{,}259.71$.}

\prob How much extra does compound interest earn in Problem 22? \answerBlank[3cm]
\answerExplain{$\$19.71$ more}{$1{,}259.71 - 1{,}240 = \$19.71$ more with compound interest.}
\end{practiceBox}

% ── True or False ────────────────────────────────────────────────────────────
\begin{practiceBox}[funPurple]{True or False?}

\trueOrFalse{Compound interest always earns more than simple interest over the same period.}
\answerExplain{True}{Compound interest earns interest on previously earned interest, so the total is always higher than simple interest (for $t > 1$ year).}

\trueOrFalse{If you deposit $\$100$ at $5\%$ compounded yearly for $2$ years, you earn exactly $\$10$ in interest.}
\answerExplain{False}{Simple interest would give $\$10$. Compound gives: Year 1: $\$5$; Year 2: $105 \times 0.05 = \$5.25$. Total interest $= \$10.25$, not $\$10$.}
\end{practiceBox}

% ── Word Problem ─────────────────────────────────────────────────────────────
\begin{practiceBox}[funBlue]{\faGlobe~Word Problems}

\wordProblem{You put $\$2{,}000$ in a savings account at $5\%$ compounded yearly. Your friend puts $\$2{,}000$ at $5\%$ simple interest. After $3$ years, who has more and by how much?}{}
\answerExplain{You have $\$15.25$ more}{Compound: $A = 2{,}000(1.05)^3 = 2{,}000 \times 1.157625 = \$2{,}315.25$. Simple: $I = 2{,}000 \times 0.05 \times 3 = \$300$; total $= \$2{,}300$. Difference: $2{,}315.25 - 2{,}300 = \$15.25$.}
\end{practiceBox}

% ── Challenge ────────────────────────────────────────────────────────────────
\begin{challengeBox}
\prob A $\$600$ loan at $8\%$ compounded yearly is not paid for $3$ years. How much is owed? How much of that is interest? \answerBlank[3cm]
\answerExplain{$\$755.83$ owed; $\$155.83$ is interest}{$A = 600(1.08)^3 = 600 \times 1.259712 = \$755.83$. Interest $= 755.83 - 600 = \$155.83$.}

\prob Explain in your own words why compound interest is sometimes called ``interest on interest.'' \answerBlank[5cm]
\answerExplain{Each year interest is added to the balance, so the next year's interest is calculated on a larger amount}{With compound interest, the interest earned in Year 1 becomes part of the principal for Year 2. So in Year 2 you earn interest on both the original amount and the Year 1 interest --- that is interest earning interest.}
\end{challengeBox}

\encouragement{You now understand how compound interest makes money grow faster over time --- a powerful concept for your financial future!}
